\chapter{図形への応用}
    \section{ある超平面から距離\texorpdfstring{$d$}{d}だけ離れた超平面の方程式}
        $n\in\naturalNumbers,\;\bm{w}\in\realNumbers^n\neq\bm{0},\;c\in\realNumbers$とする。
        写像$f:\bm{x}\in\realNumbers^n\mapsto \bm{w}^\top\bm{x} + c \in \realNumbers$が定める$\realNumbers^n$上の超平面$f(\bm{x})=0$
        から距離$d>0$だけ離れた2枚の超平面を求める方法を述べる。
        \par
        新しい写像$g = f/\norm{\bm{w}}_2$を考えると、$\nabla g = \bm{w}/\norm{w}_2,\;\norm{\nabla g}_2 = 1$であるので、超平面$g(\bm{x})=0$(これは超平面$f(\bm{x})=0$と一致)からそれと直角な方向、すなわち$\pm\nabla g$に距離$d$だけ進んだ位置では$g$の値が$\pm d$だけ変化する。
        逆に$g$の値が$\pm d$である位置は超平面$g(\bm{x})=0$から$\pm\nabla g$に距離$d$だけ離れた位置である。
        よって求めたい2平面の方程式は次式で表される。
        \[ g(\bm{x}) \pm d = 0 \iff f(\bm{x}) \pm d\norm{w}_2 = 0 \]
        \par
        $n=2, \bm{w} = [1,1]^\top, c=0.5, d=0.5$の例を\url{\currfiledir/mathematica/平面上の直線から距離dだけ離れた直線.nb}に示した。
    \section{\texorpdfstring{$n$}{n}次元単位球の体積\texorpdfstring{$V_n$}{Vn}は\texorpdfstring{$\pi^{n/2}(\GammaFunc{1+n/2})^{-1}$}{π\string^(n/2) Γ(1+n/2)\string^(-1)}}
        \begin{proof}
            \quad\par
            $\Omega(n,r)\;(r\in\realNumbers) \coloneq \setComprehension*{\bm{x}\in\realNumbers^n}{\norm{\bm{x}}_{2}\leq r}$とすると
            \begin{align*}
                V_n &= \integrate{\Omega(n,1)}{}{1}{}{\bm{x}} = \integrate{-1}{1}{\left(\integrate{\Omega(n-1,1-{x_n}^2)}{}{1}{}{\bm{x}}\right)}{}{x_n} = \integrate{-1}{1}{(\sqrt{1-{x_n}^2})^{n-1}V_{n-1}}{}{x_n} \\
                    &= V_{n-1}\integrate{-1}{1}{(1-x^2)^{(n-1)/2}}{}{x} = 2V_{n-1} \integrate{0}{1}{(1-x^2)^{(n-1)/2}}{}{x} \\
                    &= V_{n-1}\integrate{0}{1}{y^{-1/2}(1-y)^{(n-1)/2}}{}{y}	\quad (y = x^2\text{とおいた}) \\
                    &= V_{n-1}\Beta{\frac{1}{2}}{\frac{n+1}{2}} = V_{n-1}\frac{\GammaFunc{\frac{1}{2}}\GammaFunc{\frac{n+1}{2}}}{\GammaFunc{1+\frac{n}{2}}} \\
                    &= V_{n-1}\frac{\sqrt{\pi}\GammaFunc{\frac{n+1}{2}}}{\GammaFunc{\frac{(n+1)+1}{2}}} \\
                \therefore\; V_n &= V_1\prod_{i=3}^{n+1}{\sqrt{\pi}\frac{\GammaFunc{i/2}}{\GammaFunc{(i+1)/2}}} = 2\pi^{(n-1)/2}\frac{\GammaFunc{3/2}}{\GammaFunc{1+n/2}} = \frac{\pi^{n/2}}{\GammaFunc{1+n/2}}
            \end{align*}
        \end{proof}